% !PRINTED 430
\weberpage{430}

In der Kette der Kompositionen
\weq{3}{P,\ P^2,\ P^3,\ \ldots}
kann eine Form vorkommen, die in die \gesperrt{Hauptklasse} gehört, die also mit $(1,\,0,\,\tfrac14 m)$ oder $[1,\,1,\,\tfrac14(m+1)]$ äquivalent ist. Wenn dies bei $P^{\varepsilon}$ eintritt, so ist $p^{\varepsilon}$ durch die Hauptform darstellbar, und es gibt eine \gesperrt{eigentliche} Darstellung
\weq{4}{4\,p^{\varepsilon} = y^2 - D\,x^2.}
Ist umgekehrt eine solche Darstellung möglich, so gibt es in der Hauptklasse eine Form $(p^{\varepsilon},\,b,\,c)$.

\gesperrt{Ist $\varepsilon$ der kleinste positive Exponent, für den die Gleichung (4) eigentlich lösbar ist, so heißt $\varepsilon$ der Index von $p$, und $p$ gehört zum Exponenten $\varepsilon$.}

Ist
\weq{5}{P = (p,\,b,\,p^{\varepsilon-1}c),}
so ist
\[
b^2 \equiv D \pmod{4\,p^{\varepsilon}},
\]
und wenn diese Bedingung erfüllt ist, so ist in der Form $(p,\,b,\,c')$ der Diskriminante $D$ der dritte Koeffizient durch $p^{\varepsilon-1}$ teilbar. Die Form (5) bleibt erhalten, wenn wir $b$ durch irgend eine nach dem Modul $2\,p^{\varepsilon}$ kongruente Zahl $b'$ ersetzen. Ist also $a$ nicht durch $p$ teilbar und
\[
\varphi = (a,\,B,\,C)
\]
irgend eine mit $P$ einhellige Form der Diskriminante $D$, so lassen sich die beiden Kongruenzen
\[
x \equiv b \pmod{2\,p^{\varepsilon}}, \quad x \equiv B \pmod{2\,a}
\]
zugleich befriedigen, und wenn man dieses $x$ also an Stelle von $b$ und $B$ setzt, folgt:

III. \gesperrt{Ist $l$ die Klasse von $P$ und $k$ eine beliebige Klasse der Diskriminante $D$, so kann man in $l$ und $k$ die Repräsentanten wählen:}
\weq{6}{(p,\,b,\,a\,c\,p^{\varepsilon-1}),\quad (a,\,b,\,c\,p^{\varepsilon}),}
und in der komponierten Klasse $lk$ erhält man einen Repräsentanten
\weq{7}{(a\,p,\,b,\,p^{\varepsilon-1}c).}
Ist also
\[
\omega = \frac{-b+\sqrt{-m}}{2\,a}
\]
eine Wurzel der Klasse $k$, so ist $\omega\colon p$ eine Wurzel der Klasse $lk$.
